% GENERATED FILE. Edit canonical lessons, then run npm run generate:books.
% canonical-source-hash: fnv1a64:36063a8ae9597cfc
% canonical-lessons: HI-C347-paks, HI-C347-khilaf, HI-C347-virodhkarna, HI-C347-samarthankarna, HI-C347-sahmathona

\chapter{Taking Sides}
\label{ch:hi-taking-sides}
\hlchaptermodality{347}

\begin{chapteropening}
\textbf{By the end of this chapter:} \emph{I can say who is for an idea and who is against it.}

Complete the last lesson of chapter 347: I can say who is for an idea and who is against it.
\end{chapteropening}

\section[\texorpdfstring{\hi{पक्ष} (pakṣ)}{पक्ष (pakṣ)}]{\hi{पक्ष} (pakṣ) — a side (in a dispute)}

\label{lesson:HI-C347-paks}



\begin{quote}
Before the new one: say the Hindi for to interrupt, then the Hindi for but rather.
\end{quote}

\subsection*{You'll want to know: \hi{पक्ष}}
\textbf{\hi{पक्ष}} — \emph{pakṣ} — \textquotedblleft{}a side (in a dispute)\textquotedblright{}.

\begin{quote}
\hi{बहस} \hi{में} \hi{मैं} \hi{मीरा} \hi{के} \hi{पक्ष} \hi{में} \hi{हूँ।}
\end{quote}

\begin{quote}
\emph{bahas meṁ maiṁ Mīrā ke pakṣ meṁ hūṁ.}
\end{quote}

\begin{quote}
In the argument I'm on Mira's side.
\end{quote}

Masculine.

\begin{cousinweb}
\textbf{\hi{पक्ष}} (\emph{pakṣ}) is a tatsama, borrowed from Sanskrit \emph{pakṣa}, \textquotedblleft{}a wing, a side\textquotedblright{} (Monier-Williams).
\end{cousinweb}

\begin{grammarlens}[title={What you've built}]
One more word for giving reasons and opinions.
\end{grammarlens}

\subsection*{Your turn}
\begin{itemize}
\raggedright
  \item \emph{Say it:} \emph{pakṣ}
  \item \emph{Say it:} \emph{pakṣ}, once more
  \item \emph{Say it:} \emph{bahas meṁ maiṁ Mīrā ke pakṣ meṁ hūṁ.} — In the argument I'm on Mira's side.
  \item \emph{Recall:} say the Hindi for to interrupt, then the Hindi for but rather, then say \emph{pakṣ} again
\end{itemize}

\begin{tcolorbox}[breakable,colback=teal!4,colframe=teal!35!black,title={Before you move on}]
What does \hi{पक्ष} mean? (\textquotedblleft{}A side (in a dispute)\textquotedblright{}.) Say it once more.
\end{tcolorbox}
\section[\texorpdfstring{\hi{ख़िलाफ़} (khilāf)}{ख़िलाफ़ (khilāf)}]{\hi{ख़िलाफ़} (khilāf) — against (\hi{के} \hi{ख़िलाफ़})}

\label{lesson:HI-C347-khilaf}



\begin{quote}
Before the new one: say the Hindi for but rather, then the Hindi for a side.
\end{quote}

\subsection*{You'll want to know: \hi{ख़िलाफ़}}
\textbf{\hi{ख़िलाफ़}} — \emph{khilāf} — \textquotedblleft{}against (\hi{के} \hi{ख़िलाफ़})\textquotedblright{}.

\begin{quote}
\hi{गाँव} \hi{के} \hi{लोग} \hi{नई} \hi{सड़क} \hi{के} \hi{ख़िलाफ़} \hi{हैं।}
\end{quote}

\begin{quote}
\emph{gā̃v ke log naī saṛak ke khilāf haiṁ.}
\end{quote}

\begin{quote}
The villagers are against the new road.
\end{quote}

\begin{cousinweb}
\textbf{\hi{ख़िलाफ़}} (\emph{khilāf}) is usually traced to Arabic \emph{khilāf}, \textquotedblleft{}opposition, the contrary\textquotedblright{}, which reached Hindi through Persian.
\end{cousinweb}

\begin{grammarlens}[title={What you've built}]
One more word for giving reasons and opinions.
\end{grammarlens}

\subsection*{Your turn}
\begin{itemize}
\raggedright
  \item \emph{Say it:} \emph{khilāf}
  \item \emph{Say it:} \emph{khilāf}, once more
  \item \emph{Say it:} \emph{gā̃v ke log naī saṛak ke khilāf haiṁ.} — The villagers are against the new road.
  \item \emph{Recall:} say the Hindi for but rather, then the Hindi for a side, then say \emph{khilāf} again
\end{itemize}

\begin{tcolorbox}[breakable,colback=teal!4,colframe=teal!35!black,title={Before you move on}]
What does \hi{ख़िलाफ़} mean? (\textquotedblleft{}Against (\hi{के} \hi{ख़िलाफ़})\textquotedblright{}.) Say it once more.
\end{tcolorbox}
\section[\texorpdfstring{\hi{विरोध} \hi{करना} (virodh karnā)}{विरोध करना (virodh karnā)}]{\hi{विरोध} \hi{करना} (virodh karnā) — to oppose}

\label{lesson:HI-C347-virodhkarna}



\begin{quote}
Before the new one: say the Hindi for a side, then the Hindi for against.
\end{quote}

\subsection*{You'll want to know: \hi{विरोध} \hi{करना}}
\textbf{\hi{विरोध} \hi{करना}} — \emph{virodh karnā} — \textquotedblleft{}to oppose\textquotedblright{}.

\begin{quote}
\hi{किसान} \hi{नए} \hi{क़ानून} \hi{का} \hi{विरोध} \hi{कर} \hi{रहे} \hi{हैं।}
\end{quote}

\begin{quote}
\emph{kisān nae qānūn kā virodh kar rahe haiṁ.}
\end{quote}

\begin{quote}
The farmers are opposing the new law.
\end{quote}

Past, with \hi{ने} (he or she …): \hi{उसने} \hi{क़ानून} \hi{का} \hi{विरोध} \hi{किया}, \emph{usne qānūn kā virodh kiyā}.

\begin{cousinweb}
\textbf{\hi{विरोध} \hi{करना}} (\emph{virodh karnā}) joins \hi{विरोध} to \hi{करना}, \textquotedblleft{}to do\textquotedblright{}. \hi{विरोध} is a tatsama, borrowed from Sanskrit \emph{virodha}, \textquotedblleft{}opposition, hostility\textquotedblright{} (Monier-Williams).
\end{cousinweb}

\begin{grammarlens}[title={What you've built}]
One more word for giving reasons and opinions.
\end{grammarlens}

\subsection*{Your turn}
\begin{itemize}
\raggedright
  \item \emph{Say it:} \emph{virodh karnā}
  \item \emph{Say it:} \emph{virodh karnā}, once more
  \item \emph{Say it:} \emph{kisān nae qānūn kā virodh kar rahe haiṁ.} — The farmers are opposing the new law.
  \item \emph{Recall:} say the Hindi for a side, then the Hindi for against, then say \emph{virodh karnā} again
\end{itemize}

\begin{tcolorbox}[breakable,colback=teal!4,colframe=teal!35!black,title={Before you move on}]
What does \hi{विरोध} \hi{करना} mean? (\textquotedblleft{}To oppose\textquotedblright{}.) Say it once more.
\end{tcolorbox}
\section[\texorpdfstring{\hi{समर्थन} \hi{करना} (samarthan karnā)}{समर्थन करना (samarthan karnā)}]{\hi{समर्थन} \hi{करना} (samarthan karnā) — to support (an idea, a person)}

\label{lesson:HI-C347-samarthankarna}



\begin{quote}
Before the new one: say the Hindi for against, then the Hindi for to oppose.
\end{quote}

\subsection*{You'll want to know: \hi{समर्थन} \hi{करना}}
\textbf{\hi{समर्थन} \hi{करना}} — \emph{samarthan karnā} — \textquotedblleft{}to support (an idea, a person)\textquotedblright{}.

\begin{quote}
\hi{मैं} \hi{तुम्हारे} \hi{विचार} \hi{का} \hi{समर्थन} \hi{करता} \hi{हूँ।}
\end{quote}

\begin{quote}
\emph{maiṁ tumhāre vicār kā samarthan kartā hūṁ.}
\end{quote}

\begin{quote}
I support your idea.
\end{quote}

Past, with \hi{ने} (he or she …): \hi{उसने} \hi{अरुण} \hi{का} \hi{समर्थन} \hi{किया}, \emph{usne Aruṇ kā samarthan kiyā}.

\begin{cousinweb}
\textbf{\hi{समर्थन} \hi{करना}} (\emph{samarthan karnā}) joins \hi{समर्थन} to \hi{करना}, \textquotedblleft{}to do\textquotedblright{}. \hi{समर्थन} is a tatsama, borrowed from Sanskrit \emph{samarthana}, \textquotedblleft{}upholding; reconciling\textquotedblright{} (Monier-Williams).
\end{cousinweb}

\begin{grammarlens}[title={What you've built}]
One more word for giving reasons and opinions.
\end{grammarlens}

\subsection*{Your turn}
\begin{itemize}
\raggedright
  \item \emph{Say it:} \emph{samarthan karnā}
  \item \emph{Say it:} \emph{samarthan karnā}, once more
  \item \emph{Say it:} \emph{maiṁ tumhāre vicār kā samarthan kartā hūṁ.} — I support your idea.
  \item \emph{Recall:} say the Hindi for against, then the Hindi for to oppose, then say \emph{samarthan karnā} again
\end{itemize}

\begin{tcolorbox}[breakable,colback=teal!4,colframe=teal!35!black,title={Before you move on}]
What does \hi{समर्थन} \hi{करना} mean? (\textquotedblleft{}To support (an idea, a person)\textquotedblright{}.) Say it once more.
\end{tcolorbox}
\section[\texorpdfstring{\hi{सहमत} \hi{होना} (sahmat honā)}{सहमत होना (sahmat honā)}]{\hi{सहमत} \hi{होना} (sahmat honā) — to agree}

\label{lesson:HI-C347-sahmathona}



\begin{quote}
Before the new one: say the Hindi for to oppose, then the Hindi for to support.
\end{quote}

\subsection*{You'll want to know: \hi{सहमत} \hi{होना}}
\textbf{\hi{सहमत} \hi{होना}} — \emph{sahmat honā} — \textquotedblleft{}to agree\textquotedblright{}.

\begin{quote}
\hi{मैं} \hi{तुमसे} \hi{सहमत} \hi{हूँ}, \hi{लेकिन} \hi{मीरा} \hi{सहमत} \hi{नहीं} \hi{है।}
\end{quote}

\begin{quote}
\emph{maiṁ tumse sahmat hūṁ, lekin Mīrā sahmat nahīṃ hai.}
\end{quote}

\begin{quote}
I agree with you, but Mira doesn't.
\end{quote}

Past (he …): \hi{वह} \hi{सहमत} \hi{हुआ}, \emph{vah sahmat huā}.

\begin{cousinweb}
\textbf{\hi{सहमत} \hi{होना}} (\emph{sahmat honā}) joins \hi{सहमत} to \hi{होना}, \textquotedblleft{}to be, to become\textquotedblright{}. \hi{सहमत} is built from two Sanskrit words, \emph{saha}, \textquotedblleft{}together with\textquotedblright{}, and \emph{mata}, \textquotedblleft{}an opinion\textquotedblright{} (Monier-Williams): of one opinion with someone.
\end{cousinweb}

\begin{grammarlens}[title={What you've built}]
One more word for giving reasons and opinions.
\end{grammarlens}

\subsection*{Your turn}
\begin{itemize}
\raggedright
  \item \emph{Say it:} \emph{sahmat honā}
  \item \emph{Say it:} \emph{sahmat honā}, once more
  \item \emph{Say it:} \emph{maiṁ tumse sahmat hūṁ, lekin Mīrā sahmat nahīṃ hai.} — I agree with you, but Mira doesn't.
  \item \emph{Recall:} say the Hindi for to oppose, then the Hindi for to support, then say \emph{sahmat honā} again
\end{itemize}

\begin{tcolorbox}[breakable,colback=teal!4,colframe=teal!35!black,title={Before you move on}]
What does \hi{सहमत} \hi{होना} mean? (\textquotedblleft{}To agree\textquotedblright{}.) Say it once more.
\end{tcolorbox}
